% Part: computability
% Chapter: machines-computations
% Section: configuration

\documentclass[../../../include/open-logic-section]{subfiles}

\begin{document}

\olfileid{cmp}{tur}{con}
\olsection{స్థితివిన్యాసాలు, గణనలు}

\begin{explain}
ఇంతకుముందు సరి యంత్రపు స్థితివిన్యాసాలను వరుసగా అనుసరించిన సంగతి
గుర్తుచేసుకోండి. టేపు, చదువు--వ్రాత శీర్షం, ట్యూరింగ్ యంత్ర ప్రోగ్రాముతో కూడిన
ఊహాత్మక యంత్రాంగం, ట్యూరింగ్ యంత్ర గణన అంటే ఏమిటో దృశ్యీకరించడానికి ఒక
అంతర్బోధాత్మక పద్ధతి మాత్రమే. ఆచారబద్ధంగా, ఇచ్చిన ఇన్‌పుట్‌పై ట్యూరింగ్
యంత్రపు గణనను \emph{స్థితివిన్యాసాల} క్రమంగా నిర్వచించవచ్చు. ఒక
స్థితివిన్యాసమేమో, సంకేతాల క్రమం (గణనలో ఒక నిర్దిష్ట బిందువు వద్ద టేపు
విషయానికి అనుగుణంగా ఉంటుంది), చదువు--వ్రాత శీర్షం స్థానాన్ని సూచించే ఒక
సంఖ్య, ఒక స్థితితో కూడి ఉంటుంది. వీటిని ఉపయోగించి, ఇచ్చిన ఇన్‌పుట్‌పై ట్యూరింగ్
యంత్రం~$M$ ఏమి గణిస్తుందో నిర్వచించవచ్చు.
\end{explain}

\begin{defn}[స్థితివిన్యాసం]
ట్యూరింగ్ యంత్రం $M = \tuple{Q, \Sigma, q_0, \delta}$ యొక్క
\emph{స్థితివిన్యాసం} అనేది $\tuple{C, m, q}$ అనే త్రయం; ఇక్కడ
\begin{enumerate}
\item $C \in \Sigma^*$ అనేది $\Sigma$లోని సంకేతాల పరిమిత క్రమం,
\item $m \in \Nat$ అనేది $< \len{C}$ అయిన సంఖ్య, అలాగే
\item $q \in Q$.
\end{enumerate}
అంతర్బోధాత్మకంగా, క్రమం~$C$ టేపు విషయం (అత్యంత ఎడమ గడి నుండి చివరి ఖాళీ-కాని
లేదా ఇంతకుముందు సందర్శించిన గడి వరకు అన్ని గడుల సంకేతాలు); $m$ అనేది
చదువు--వ్రాత శీర్షం చదువుతున్న గడి సంఖ్య (అత్యంత ఎడమ గడి సంఖ్య $0$తో
మొదలవుతుంది); $q$ యంత్రపు ప్రస్తుత స్థితి.
\end{defn}

\begin{explain}
ట్యూరింగ్ యంత్రానికి సాధ్యమైన ఇన్‌పుట్ ఒక సంకేతాల క్రమం; సాధారణంగా అది ఏదో
ఒక రూపంలో సంఖ్యను సంకేతీకరించే క్రమం. ఆ ఇన్‌పుట్‌పై పనిచేయడానికి ట్యూరింగ్
యంత్రాన్ని ప్రారంభించే స్థితివిన్యాసమే దాని ప్రారంభ స్థితివిన్యాసం: టేపుపై
టేపు-అంత గుర్తు ఉండి, దాని వెంటనే కుడివైపు గడులపై వ్రాసిన ఇన్‌పుట్ ఉంటుంది;
చదువు--వ్రాత శీర్షం ఇన్‌పుట్ యొక్క అత్యంత ఎడమ గడిని (అంటే ఎడమ చివరి గుర్తుకు
కుడివైపున ఉన్న గడిని) చదువుతుంది; యంత్రాంగం నిర్దేశిత ప్రారంభ స్థితి~$q_0$లో
ఉంటుంది.
\end{explain}

\begin{defn}[ప్రారంభ స్థితివిన్యాసం]
ఇన్‌పుట్ $I \in \Sigma^*$కు $M$ యొక్క \emph{ప్రారంభ స్థితివిన్యాసం}
\[
\tuple{\TMendtape \frown I, 1, q_0}.
\]
\end{defn}

\begin{explain}
$\frown$ సంకేతం \emph{అనుసంధానాన్ని} సూచిస్తుంది---ఇన్‌పుట్ సంకేతమాల ఎడమ
చివరి గుర్తుకు వెంటనే కుడివైపున మొదలవుతుంది.
\sourcecorrection{OLTETURMAC-003}{ఆంగ్ల మూలం ఇన్‌పుట్ ``begins immediately to the left end marker'' అని అసంపూర్ణంగా రాసింది. నిర్వచనంలోని $\TMendtape \frown I$కూ ముందరి వివరణకూ అనుగుణంగా, ఇన్‌పుట్ ఎడమ చివరి గుర్తుకు వెంటనే కుడివైపున మొదలవుతుందని స్పష్టం చేశాం.}
\end{explain}

\begin{defn}
(యంత్రం~$M$ ప్రకారం) స్థితివిన్యాసం $\tuple{C, m, q}$
\emph{ఒక అంచెలో స్థితివిన్యాసం $\tuple{C', m', q'}$ను ఇస్తుంది} అని
చెప్పేది, అప్పుడే మాత్రమే:
\begin{enumerate}
\item $C$ యొక్క $m$వ సంకేతం $\sigma$,
\item $M$ యొక్క నిర్దేశ సమితి $\delta(q, \sigma) =
  \tuple{q', \sigma', D}$ అని నిర్దేశిస్తుంది,
\item $C'$ యొక్క $m$వ సంకేతం $\sigma'$, అలాగే
\item
\begin{enumerate}
\item $D = L$, $m>0$ అయితే $m' = m - 1$; లేకపోతే $m'=0$, లేదా
\item $D = R$, $m' = m + 1$, లేదా
\item $D = N$, $m' = m$,
\end{enumerate}
\item $m' = \len{C}$ అయితే, $\len{C'} = \len{C} + 1$, అలాగే $C'$ యొక్క
  $m'$వ సంకేతం~$\TMblank$. లేకపోతే $\len{C'}=\len{C}$.
\item $i < \len{C}$, $i \neq m$ అయ్యే ప్రతి $i$కు $C'(i) = C(i)$.
\end{enumerate}
\end{defn}

\begin{defn}\ollabel{defn:run-output}
\emph{ఇన్‌పుట్~$I$పై $M$ యొక్క నడక} అనేది $M$ యొక్క పరిమిత లేదా అనంత
స్థితివిన్యాసాల క్రమం $C_i$; ఇక్కడ $C_0$ అనేది ఇన్‌పుట్~$I$కు $M$ యొక్క
ప్రారంభ స్థితివిన్యాసం, క్రమంలో $C_{i+1}$ ఉన్న ప్రతి $i$కు $C_i$ ఒక అంచెలో
$C_{i+1}$ను ఇస్తుంది.
\sourcecorrection{OLTETURMAC-004}{ఆంగ్ల మూలం నడకలో ప్రతి $C_i$ తదుపరి $C_{i+1}$ను ఇస్తుందని షరతు విధించి, వెంటనే చివరి స్థితివిన్యాసం వద్ద నిర్వచితం కాని సంక్రమణతో నడక ఆగవచ్చని చెప్పింది. ఆగే నడకల ఉనికిని కాపాడటానికి నడక పరిమితం లేదా అనంతం కావచ్చనీ, క్రమంలో వారసుడు ఉన్న సూచికలకే ఒక-అంచె షరతు వర్తిస్తుందనీ స్పష్టం చేశాం.}

$C_k = \tuple{C, m, q}$ అయి, $C$ యొక్క $m$వ సంకేతం~$\sigma$ అయి,
$\delta(q, \sigma)$ నిర్వచితం కాకపోతే, $M$ \emph{ఇన్‌పుట్ $I$పై $k$
అంచెల తరువాత ఆగుతుంది} అంటాం. ఆ సందర్భంలో, ఇన్‌పుట్~$I$కు $M$ యొక్క
\emph{అవుట్‌పుట్}~$O$; ఇక్కడ $O$ అనేది $\TMblank$తో ముగియని సంకేతమాల,
ఏదో~$j \in \Nat$కు $C = \TMendtape \concat O \concat \TMblank^j$.
($\TMblank^j$ అనేది $j$ $\TMblank$ల క్రమం.)
\end{defn}

\begin{explain}
ఈ నిర్వచనం ప్రకారం, $M$ యొక్క అవుట్‌పుట్~$O$ ఎల్లప్పుడూ $\TMblank$ కాని
సంకేతంతో ముగుస్తుంది; లేదా సమయం~$k$ వద్ద అత్యంత ఎడమ $\TMendtape$ తప్ప టేపు
మొత్తం $\TMblank$లతో నిండి ఉంటే, $O$ ఖాళీ సంకేతమాల అవుతుంది.
\end{explain}

\end{document}
